\subsection{例外集}\label{sec:5i}
在下面的整个讨论中，$\A$ 始终是 c-有限、三角化的 $\Ai$ 范畴。
\begin{definition*}
若 $\A$ 的有限有序对象族 $Y=(Y_1,\ldots,Y_m)$ 满足以下两个条件，就称它为\emph{例外集}。
\begin{itemize}
\item 每个 $Y_k$ 都是\emph{例外对象}；按定义，这意味着其自同态空间由恒等态射张成：
\[
\Hom_{H(\A)}(Y_k,Y_k)=\K\cdot[e_{Y_k}].
\]
\item 在上同调层次上，不同 $Y_k$ 之间的态射具有方向性，即
\[
\Hom_{H(\A)}(Y_k,Y_l)=0\quad(k>l).
\]
\end{itemize}
若此外 $Y_k$ 在\secref{3j}意义下生成 $\A$，就称 $Y$ 为\emph{满例外集}。\znote{76}{“例外对象”一词来自 Bondal 的工作：在三角范畴中，若对象 $E$ 满足 $\Hom(E,E)=\K$ 且 $\Hom(E,E[i])=0$（$i\ne0$），就称 $E$ 例外；两两例外、且态射只有单向的一族例外对象称为例外集，若它还生成整个范畴，就称为满的。对满例外集，任意对象都可沿这族对象作有限阶梯式的 Postnikov 分解，这正是引理 5.13 与 (5.23) 的内容。[来源：nLab "exceptional collection"，https:/\allowbreak /\allowbreak ncatlab.org/\allowbreak nlab/\allowbreak show/\allowbreak exceptional+\allowbreak collection（2026-11-15 访问）。]}
\end{definition*}
\begin{lemma}\label{st:5.13}
设 $Y$ 为满例外集。那么，对任意 $X$，在 $H^0(\A)$ 中有同构
\[
T_{Y_1}\cdots T_{Y_m}(X)\cong0.
\]
\end{lemma}
\SourcePage{84}
\begin{proof}
如\secref{5e}，令 $X_m=X$、$X_{k-1}=T_{Y_k}(X_k)$。目标是证明，对所有 $X$ 都有 $X_0\cong0$。先考虑 $X=Y_l$ 的情形。只要 $k>l$，就有 $\Hom_{H(\A)}(Y_k,X)=0$，因此由\eref{5.4}，$X_{k-1}\cong X_k\cong\cdots\cong X$。下一步中，三角化简为
\SourceDiagram{diagram-5-13}{0.27830739}
因此 $X_{l-1}\cong\Cone(e_{Y_l})\cong0$。此后，所有较低的 $X_k$（$k<l$）也立即可见为零。

$\F=T_{Y_1}\cdots T_{Y_m}$ 是 $\A$ 到自身的 $\Ai$ 函子，因此由\stref{推论}{3.9}保持正合三角。这意味着，若两个对象 $X_0,X_1$ 的像 $\F(X_0),\F(X_1)$ 均为零，那么它们之间任意态射的映射锥 $X_2=\Cone(X_0\to X_1)$ 的像也为零。此外，$\F$ 在同构意义下也与平移相容，这更为显然。由于 $Y_k$ 生成 $\A$，所需结论随之成立。
\end{proof}
这个论证实际上给出了任意对象 $X$ 到 $Y_k$ 的若干平移的显式 \emph{Postnikov 分解}（在一般三角范畴中的例外集框架下，这是众所周知的，参见\bcite{66}§2）。具体地，仍令 $Z_k=\Hom_{H(\A)}(Y_k,X_k)$，有如下阶梯图，其左侧为正合三角，右侧为交换三角：
\SourceEquationDiagram{diagram-5-23}{0.28246123}{5.23}
\SourcePage{85}
\begin{note}\label{st:5.14}
由\stref{推论}{5.7}可知，若 $\A$ 有满例外集，则它分裂闭。事实上，还有稍强的结论：若 $\Q^{\mathrm{fin}}\subset\Q=\modcat(\A)$ 为如下全子范畴，其对象是对每个 $X\in\Ob\A$ 都有 $H(\M(X))$ 有限维的 $\Ai$ 模 $\M$，那么 Yoneda 嵌入给出拟等价 $\A\hookrightarrow\Q^{\mathrm{fin}}$。经典框架下的相应定理是\bcite{27}第 530 页的推论：若 $\mathcal T$ 是具有满例外集的三角范畴，则从 $\mathcal T$ 到有限维 $\K$ 向量空间范畴的任意上同调函子都可表示（受三角范畴理论的局限，这里的证明比我们所处理情形的证明困难得多）。
\end{note}

\subsection{变异}\label{sec:5j}
可以通过以下变换修改给定例外集 $Y=(Y_1,\ldots,Y_m)$：
\begin{itemize}
\item 将每个 $Y_k$ 替换为在 $H^0(\A)$ 中与之同构的对象 $\widetilde Y_k$，形成所得对象族 $\widetilde Y$。
\item 给定 $(\sigma_1,\ldots,\sigma_m)\in\Z^m$，取 $S^{\sigma}Y=(S^{\sigma_1}Y_1,\ldots,S^{\sigma_m}Y_m)$，其中 $S$ 为 $\A$ 上的平移函子。
\item 最后，第 $k$ 个\emph{基本变异}（$1\leq k<m$）为运算 $Y\mapsto U=\mathscr L_kY$，其中 $j\ne k,k+1$ 时 $U_j=Y_j$，其余两个对象为
\begin{equation}\label{eq:5.24}
U_k=T_{Y_k}Y_{k+1},\qquad U_{k+1}=Y_k.
\end{equation}
\end{itemize}
暂且在例外集的同构类层次上讨论，即把第一类变换视为平凡变换。那么，第二类与第三类变换都可逆。第三类变换的逆为 $Y\mapsto V=\mathscr L_k^{-1}Y$，其中仍在 $j\ne k,k+1$ 时取 $V_j=Y_j$，而
\[
V_k=Y_{k+1},\qquad V_{k+1}=T_{Y_{k+1}}^{\vee}Y_k.
\]
此外，还有辫关系
\[
\begin{aligned}
\mathscr L_j\mathscr L_kY&\cong\mathscr L_k\mathscr L_jY&&(|j-k|\geq2),\\
\mathscr L_j\mathscr L_{j+1}\mathscr L_jY&\cong\mathscr L_{j+1}\mathscr L_j\mathscr L_{j+1}Y.
\end{aligned}
\]
由此得到群 $\Z^m\rtimes\operatorname{Br}_m$ 在长度为 $m$ 的例外集的同构类集合上的作用。子群 $\Z^m$ 通过平移作用；辫群 $\operatorname{Br}_m$ 的标准生成元 $\beta_1,\ldots,\beta_{m-1}$ 通过基本变异作用；半直积则通过对称群 $\operatorname{Sym}_m$ 的置换表示形成，这里把它视为 $\operatorname{Br}_m$ 的商群。此外，若一个例外集是满的，那么其 $\Z^m\rtimes\operatorname{Br}_m$ 轨道上的所有例外集也都是满的。\znote{78}{平移给出 $\Z^m$，基本变异 $\mathscr L_k$ 给出辫群 $\operatorname{Br}_m$ 的标准生成元，上两式正是辫关系的两种情形（相距较远的生成元交换、相邻生成元满足辫关系）。这套作用在 Bondal 关于螺旋与 Koszul 代数的工作中已经出现：螺旋由反复作变异得到，而变异不改变导出范畴；稍后 §5n 会把这一点重述为有向 $\Ai$ 范畴上的作用。[来源：nLab "exceptional collection"，https:/\allowbreak /\allowbreak ncatlab.org/\allowbreak nlab/\allowbreak show/\allowbreak exceptional+\allowbreak collection（2026-11-15 访问）。]}
\SourcePage{86}
证明只需对正合三角作简单论证，可以在三角范畴框架中完成，而无须涉及 $\Ai$ 结构；例如见\bcite{66}命题 2.4.2。需要强调，$\mathscr L_k$ 并非函子：它们作用于 $\A$ 内的对象族，却不作用于这些对象之间的态射。

\subsection{Koszul 对偶集}\label{sec:5k}
设 $Y=(Y_1,\ldots,Y_m)$ 是 $\A$ 中的满例外集。对于 $k\leq l$，记
\[
Y_{k,l}=T_{Y_k}T_{Y_{k+1}}\cdots T_{Y_{l-1}}(Y_l),
\]
最简单的情形为 $Y_{k,k}=Y_k$。
\begin{definition*}
\emph{右 Koszul 对偶集} $Y^!$ 定义为由对象 $Y^!_{m+1-k}=Y_{1,k}$（$1\leq k\leq m$）组成的对象族。
\end{definition*}
很容易看出，它仍是满例外集，并由原来的 $Y$ 沿“整体半扭转”
\begin{equation}\label{eq:5.25}
\Delta^{1/2}=\beta_{m-1}(\beta_{m-2}\beta_{m-1})\cdots(\beta_2\cdots\beta_{m-1})(\beta_1\cdots\beta_{m-1})\in\operatorname{Br}_m
\end{equation}
变异而得（辫群中的复合从右到左进行，因此\eref{5.25}对应的变异从 $\mathscr L_{m-1}$ 开始，接着依次施加 $\mathscr L_{m-2}$ 直到 $\mathscr L_1$，如此继续）。复合有关正合三角\eref{5.4}中的映射 $[i]$，便对所有 $j<k\leq l$ 得到 $H^0(\A)$ 中的典范态射
\[
[i_{j,k,l}]:Y_{k,l}\longrightarrow Y_{k-1,l}\longrightarrow\cdots\longrightarrow Y_{j,l}.
\]
令 $[i_{k,k,l}]$ 为恒等态射，就将其延伸到 $j=k$ 的情形。特别地，对每个 $k$，有典范态射 $[i_{1,k,k}]:Y_k\to Y_{1,k}=Y^!_{m+1-k}$。现在讨论 Koszul 对偶集的一些性质。这些性质基本上都来自三角\eref{5.4}诱导的长正合列。为简洁起见，用 $\Hom$ 代替 $\Hom_{H(\A)}$；同时省略通常括在上同调层次态射外面的方括号，即以 $i_{j,k,l}$ 表示 $[i_{j,k,l}]$。
\begin{lemma}\label{st:5.15}
\textup{(i)} 当 $k=l$ 时，$\Hom(Y_l,Y^!_{m+1-k})$ 为一维，由 $i_{1,k,k}$ 生成；当 $k\ne l$ 时为零。

\textup{(ii)} $\Hom(X,Y^!_{m+1-k})\cong\Hom(Y_k,X_k)^{\vee}$，其中如\eref{5.23}，$X_k=T_{Y_{k+1}}\cdots T_{Y_m}(X)$。

\textup{(iii)} $\Hom(Y^!_{m-k},Y^!_{m+1-k})\cong\Hom(Y_k,Y_{k+1})^{\vee}[-1]$。
\end{lemma}
\begin{proof}
省略直接可得的 (i)。为由它推出更一般的 (ii)，从两端遍历分解\eref{5.23}，并施加 $\Hom(-,Y^!_{m+1-k})$，得到
\[
\begin{aligned}
\Hom(X,Y^!_{m+1-k})&\cong\Hom(X_{m-1},Y^!_{m+1-k})\cong\cdots\cong\Hom(X_k,Y^!_{m+1-k}),\\
0&\cong\Hom(X_0,Y^!_{m+1-k})\cong\cdots\cong\Hom(X_{k-1},Y^!_{m+1-k}).
\end{aligned}
\]
\SourcePage{87}
余下的正合三角联系 $X_k$ 与 $X_{k-1}$，由此得到
\[
\Hom(X_k,Y^!_{m+1-k})\cong\Hom(\Hom(Y_k,X_k)\otimes Y_k,Y^!_{m+1-k})\cong\Hom(Y_k,X_k)^{\vee}.
\]
对于 (iii)，在 (ii) 中取 $X=Y^!_{m-k}$。此时 $X_k$ 是映射 $i_{1,k+1,k+1}:Y_{k+1}\to Y^!_{m-k}$ 的锥，因此 $\Hom(Y_k,X_k)\cong\Hom(Y_k,Y_{k+1})[1]$。
\end{proof}
\begin{lemma}\label{st:5.16}
取\eref{5.23}中的一次态射 $Z_k\otimes Y_k\to Z_{k+1}\otimes Y_{k+1}$，利用上面 \textup{(ii)} 的同构，将其视为映射
\begin{equation}\label{eq:5.26}
\Hom(X,Y^!_{m+1-k})^{\vee}\otimes Y_k\longrightarrow\Hom(X,Y^!_{m-k})^{\vee}\otimes Y_{k+1}.
\end{equation}
利用 \textup{(iii)}，它可以表述如下：取 $\Hom(Y_k,Y_{k+1})$ 与 $\Hom(Y^!_{m-k},Y^!_{m+1-k})[1]$ 的对偶基 $c_l,c_l^{\vee}$。对每个 $l$，先取与 $c_l^{\vee}$ 作乘法的线性映射 $\Hom(X,Y^!_{m-k})\to\Hom(X,Y^!_{m+1-k})$，将其对偶化，再与 $c_l$ 组合，就得到\eref{5.26}两侧空间之间的映射。最后对所有 $l$ 求和，即得所需结果。\znote{79}{用对偶基 $c_l,c_l^{\vee}$ 描述映射，是把两件事配对：$u_1\cdot c_l$ 把 $\Hom(X,Y^!_{m-k})$ 的元素与 $c_l\in\Hom(Y_k,Y_{k+1})$ 复合，而 $c_l^{\vee}\cdot u_0$ 把 $c_l^{\vee}$ 作用到 $\Hom(X,Y^!_{m+1-k})$ 上。两项互为转置，故求和后得到 (5.26) 两侧空间之间的恰当映射；谱序列微分 (5.29) 中的符号 $(-1)^{|c_l|(|u_0|+|u_1|)}$ 则来自把 $c_l^{\vee}$ 移过 $u_0$ 时交换两个奇次元素所生的 Koszul 符号。[来源：并见 Seidel, Fukaya categories and Picard–Lefschetz theory, §5k–5l（引理 5.15、5.16）。]}
\end{lemma}
\begin{proof}[证明概要]
先假设 $X$ 只是某个态射 $Z_k\otimes Y_k\to Z_{k+1}\otimes Y_{k+1}$ 的锥。在这一情形下，前一引理的所有同构都容易写出，可以通过显式计算验证结论（可以使用三角范畴的标准工具抽象计算，也可以用扭曲复形计算）。现在考虑一般对象 $X$。取典范映射 $X\to X_{k-1}$，将其补成正合三角 $X'\to X\to X_{k-1}$。再应用标准正合三角，可见 $X'\to X$ 在与 $X$ 相应的\eref{5.26}两侧空间及与 $X'$ 相应的对应空间之间诱导同构，且与两侧之间的映射交换。类似地，若令 $X''=X'_{k+1}$，则映射 $X'\to X''$ 也在\eref{5.26}上诱导同构。因此，只需对 $X''$ 证明交换性。但它满足 $T_{Y_k}T_{Y_{k+1}}X''=X'_{k-1}=0$，因此是上述类型的锥。
\end{proof}

\subsection{Beilinson 型谱序列}\label{sec:5l}
设对象 $X$ 具有 Postnikov 分解\eref{5.23}。令 $K$ 为上同调函子，即从 $H(\A)$ 到分次 $\K$ 向量空间的反变函子，要求它与有限直和及平移相容，并将正合三角映为长正合列；精确定义可见\bcite{61}、\bcite{143}。存在收敛到 $K(X)$ 的谱序列，其起始项为
\begin{equation}\label{eq:5.27}
E_1^{rs}=K^{r+s}(Z_{m-r}\otimes Y_{m-r})\cong(Z_{m-r}^{\vee}\otimes K(Y_{m-r}))^{r+s}.
\end{equation}
微分 $\partial_1^{r\bullet}$ 是\eref{5.23}中态射 $Z_{m-r-1}\otimes Y_{m-r-1}\to Z_{m-r}\otimes Y_{m-r}$ 在 $K$ 下的像。为建立这一谱序列，可以采用
\SourcePage{88}
Atiyah–Hirzebruch 谱序列的标准定义\bcite{138}（第 15 章），分别用 Postnikov 分解与上同调函子代替胞腔分解与广义同调理论；也可以采用正合偶的表述\bcite{61}（第 263 页）。将这一一般构造用于当前情形，得到：
\begin{proposition}\label{st:5.17}
对 $H(\A)$ 的任意一对对象 $U_0,U_1$，存在收敛到 $\Hom(U_0,U_1)$ 的谱序列，其起始项为
\begin{equation}\label{eq:5.28}
E_1^{rs}=(\Hom(U_0,Y^!_{r+1})\otimes\Hom(Y_{m-r},U_1))^{r+s}.
\end{equation}
按\stref{引理}{5.16}选取对偶基，其中 $k=m-r-1$。那么微分 $\partial^{r\bullet}:E_1^{r\bullet}\to E_1^{r+1,\bullet}$ 为
\begin{equation}\label{eq:5.29}
\partial(u_0\otimes u_1)=\sum_l(-1)^{|c_l|(|u_0|+|u_1|)}(c_l^{\vee}\cdot u_0)\otimes(u_1\cdot c_l).
\end{equation}
\end{proposition}
\begin{proof}
取 $X=U_0$、$K=\Hom(-,U_1)$。利用同构 $Z_{m-r}^{\vee}=\Hom(Y_{m-r},X_{m-r})^{\vee}\cong\Hom(X,Y^!_{r+1})$，起始项\eref{5.27}就成为\eref{5.28}，而\stref{引理}{5.16}给出公式\eref{5.29}（严格地说，符号问题还有些含糊，因为我们并未完整说明谱序列如何建立；不过此处并不需要完全精确）。\znote{80}{这一构造是把拓扑中的 Atiyah–Hirzebruch 谱序列移植到三角范畴：把空间换成对象 $X$、把胞腔分解换成 Postnikov 分解 (5.23)、把广义同调换成上同调函子 $K$，$E_1$ 项就是各层“胞腔”的上同调，微分由粘合映射诱导。经典情形是 Beilinson 对射影空间上结构层与对角线的分解，它给出 $\mathbb{P}^n$ 上相干层上同调的谱序列。[来源：并见 Seidel, Fukaya categories and Picard–Lefschetz theory, §5l 与笔记 5.18 所引文献。]}
\end{proof}
\begin{note}\label{st:5.18}
谱序列\eref{5.28}最初的例子来自 Beilinson 对射影空间上对角线的分解，见\bcite{19}或\bcite{104}第 245 页。它在\bcite{67}中被推广到凝聚层的例外集；从这里推广到一般三角范畴相当直接。完整论述见\bcite{66}，特别是其中的 2.7.3 与 3.3.2 节（注意，由于使用左 Koszul 对偶而非右对偶，他们的谱序列 $E_1$ 项与我们的略有不同）。与 Koszul 对偶性的联系在\bcite{26}、\bcite{25}中有解释；另见下面的\secref{5o}。
\end{note}

\clearpage
\subsection{有向\texorpdfstring{$\Ai$}{A∞} 范畴}\label{sec:5m}
到目前为止，我们采用了外在的观点，在一个固定的三角化 $\Ai$ 范畴内部研究例外集及其变异。另一种办法是只考虑一类特殊的 $\Ai$ 范畴：它们的对象本身在导出范畴中构成例外集。在这种内在的背景下，变异就是将这样的范畴变成另一个具有等价导出范畴的范畴的操作。
\begin{definition*}
\emph{有向 $\Ai$ 范畴}是一个具有 $m$ 个对象的严格含幺 $\Ai$ 范畴 $\A^{\to}$，连同这些对象的一个排序 $(Y_1,\ldots,Y_m)$，使
\[
\hom_{\A^{\to}}(Y_i,Y_j)=\begin{cases}
0,&i>j,\\\K\cdot e_{Y_i},&i=j,\\\K\text{ 上的有限维空间},&i<j.
\end{cases}
\]

\end{definition*}
\SourcePage{89}
由于严格含幺，只有下列复合映射不是由单位元条件直接确定的：
\begin{equation}\label{eq:5.30}
\begin{aligned}
\mu^d_{\A^{\to}}:\hom_{\A^{\to}}(Y_{i_{d-1}},Y_{i_d})\otimes\cdots\otimes\hom_{\A^{\to}}(Y_{i_0},Y_{i_1})\\
\longrightarrow\hom_{\A^{\to}}(Y_{i_0},Y_{i_d})[2-d],\qquad i_0<\cdots<i_d.
\end{aligned}
\end{equation}
所以 $\A^{\to}$ 可以由有限多个 $\K$ 中的元素完全描述，即\eref{5.30}的矩阵系数。令 $\A$ 为 $\A^{\to}$ 的三角化包络。由定义，$Y_k$ 在上同调全忠实函子 $\A^{\to}\to\A$ 下的像构成一个例外集。由构造，$\A$ 由 $\A^{\to}$ 的对象生成，因此这个例外集是满的。
\begin{note}\label{st:5.19}
更具体地，取 $\A=\Tw\A^{\to}$。有向性使扭曲复形的定义在技术上可以作若干小的简化。首先，任意这样的复形都可以写成
\begin{equation}\label{eq:5.31}
\left(X=\bigoplus_{i=1}^mV^i\otimes Y_i,\quad\delta_X=(\delta_X^{ji})\right),
\end{equation}
其中 $V^1,\ldots,V^m$ 为有限维分次向量空间，而 $\delta_X^{ji}\in\Hom_{\K}(V^i,V^j)\otimes\hom_{\A}(Y_i,Y_j)$ 为一次元素。对角项形如 $\delta_X^{ii}=\partial_{V^i}\otimes e_{Y_i}$，其中 $(V^i,\partial_{V^i})$ 是向量空间的链复形。

其次，$\delta_X$ 的下三角条件自动满足，因此在这一特殊情形下，可以从扭曲复形的定义中删去。具体地，可以先用 $\bigoplus_{i\geq k}V^i\otimes Y_i$ 对\eref{5.31}进行滤过，再用首个链复形 $(V^k,\partial_{V^k})$ 的次数滤过加以细化。这样得到 $X$ 的一个滤过，相对于它，$\delta_X$ 满足所需的下三角条件。

最后，任意扭曲复形 $X$ 都在 $H^0(\A)$ 中同构于某个微分的对角项全为零的 $\widetilde X$，即 $\delta_{\widetilde X}^{ii}=0$。事实上，Postnikov 分解\eref{5.23}说明，在同构意义下，可以从 $X_0=0$ 出发，通过一系列映射锥
\begin{equation}\label{eq:5.32}
X_k\cong\Cone(S^{-1}X_{k-1}\longrightarrow Z_k\otimes Y_k)
\end{equation}
重建 $X$。假设已有与 $X_{k-1}$ 同构的对象 $\widetilde X_{k-1}$，它是只包含 $Y_1,\ldots,Y_{k-1}$ 直和项、且微分对角项为零的扭曲复形。令 $\widetilde X_k$ 为与\eref{5.32}相应的态射 $S^{-1}\widetilde X_{k-1}\to Z_k\otimes Y_k$ 的锥。它同构于 $X_k$，并具有与 $\widetilde X_{k-1}$ 相同的性质；反复进行这一构造，最终就得到所需的 $\widetilde X=\widetilde X_m$。
\end{note}

\SourcePage{90}
\subsection{有向子范畴}\label{sec:5n}
暂时只考虑所有 $\hom$ 空间都有限维的 $\Ai$ 范畴，这可以简化一些技术细节。设 $Y=(Y_1,\ldots,Y_m)$ 是这样的范畴 $\A$ 中的对象族，未必是例外集。
\begin{definition*}
相应的\emph{有向 $\Ai$ 子范畴} $\A^{\to}=\A^{\to}(Y)$ 定义为如下唯一的有向 $\Ai$ 范畴：其对象为 $Y_k$，在 $i<j$ 时 $\hom_{\A^{\to}}(Y_i,Y_j)=\hom_{\A}(Y_i,Y_j)$，而非平凡复合映射\eref{5.30}与 $\A$ 中的对应映射相同。
\end{definition*}
这一构造具有函子性：设 $\F:\A\to\B$ 为 $\Ai$ 函子，令 $\B^{\to}$ 为与像对象族 $\F(Y)$ 相应的有向 $\Ai$ 子范畴。则存在唯一的严格含幺 $\Ai$ 函子 $\F^{\to}:\A^{\to}\to\B^{\to}$，其在对象的严格递增链上的限制等于 $\F$。若 $\F$ 上同调全忠实，则 $\F^{\to}$ 为拟同构。

需要提醒一点：若 $\A$ 严格含幺，就可以把 $\A^{\to}$ 与 $\A$ 中由 $i<j$ 时的 $\hom_{\A}(Y_i,Y_j)$ 以及各 $k$ 对应的子空间 $\K\cdot e_{Y_k}$ 组成的子范畴等同。不过，一般的 c-含幺情形下，$\A^{\to}$ 严格说来并非子范畴。对此有如下补救：
\begin{lemma}\label{st:5.20}
设 $Y$ 为 $\A$ 中的对象族，$\A^{\to}$ 为相应的有向 $\Ai$ 子范畴。那么存在在对象上为恒等映射的 $\Ai$ 函子 $\mathcal H:\A^{\to}\to\A$，对所有 $i<j$ 诱导同构
\[
\Hom_{H(\A^{\to})}(Y_i,Y_j)\cong\Hom_{H(\A)}(Y_i,Y_j).
\]
\end{lemma}
\begin{proof}
先取最小模型，再用\stref{引理}{2.1}将其严格化，得到严格含幺的 $\widetilde{\A}$ 及在对象上为恒等映射的拟等价 $\Phi:\A\to\widetilde{\A}$。令 $\widetilde{\A}^{\to}$ 为与同一对象族 $Y$ 相应的有向子范畴。严格含幺性给出包含函子 $\widetilde{\mathcal H}:\widetilde{\A}^{\to}\to\widetilde{\A}$，同时有诱导拟同构 $\Phi^{\to}:\A^{\to}\to\widetilde{\A}^{\to}$。取在对象上为恒等映射的逆向拟等价 $\Psi:\widetilde{\A}\to\A$，令 $\mathcal H=\Psi\circ\widetilde{\mathcal H}\circ\Phi^{\to}$ 即可。
\end{proof}
\begin{lemma}\label{st:5.21}
设 $Y,Y'$ 是 $\A$ 中的两个对象族，并且对每个 $j$，$Y_j$ 与 $Y'_j$ 在 $H^0(\A)$ 中同构。那么相应的有向子范畴 $\A^{\to}(Y)$、$\A^{\to}(Y')$ 拟同构，且拟同构可由保持对象顺序的严格含幺 $\Ai$ 函子给出。
\end{lemma}
\begin{proof}
令 $\B\subset\A$ 为只包含对象 $Y_j$ 的全 $\Ai$ 子范畴，$\B'$ 为只包含对象 $Y'_j$ 的相应子范畴，$\C$ 为同时包含两类对象的子范畴。由假设，$\B$ 与 $\B'$ 到 $\C$ 的包含都是拟等价。因此，利用\stref{定理}{2.9}可以构造拟等价 $\B\to\B'$。由构造很容易看出，它将每个 $Y_j$ 映到 $Y'_j$，故实际上是拟同构。再转到有向子范畴即可。
\end{proof}
\begin{lemma}\label{st:5.22}
设 $Y$ 是三角化 $\Ai$ 范畴 $\A$ 中的满例外集，$\A^{\to}$ 为相应的有向 $\Ai$ 子范畴，$\mathcal H:\A^{\to}\to\A$ 为\stref{引理}{5.20}中的
\SourcePage{91}
$\Ai$ 函子。那么存在拟等价 $\G:\Tw\A^{\to}\to\A$，其在 $\A^{\to}\subset\Tw\A^{\to}$ 上的限制，在 $H^0(\operatorname{fun}(\A^{\to},\A))$ 中同构于 $\mathcal H$。特别地，$D(\A^{\to})$ 与 $H^0(\A)$ 作为三角范畴等价。
\end{lemma}
\begin{proof}
$Y$ 是例外集这一事实意味着 $\mathcal H$ 上同调全忠实。由于 $\A$ 三角化，且由 $Y$ 的对象生成，它实际上就是 $\A^{\to}$ 的三角化包络。$\Tw\A^{\to}$ 也一样；现在应用\stref{引理}{3.20}即可。
\end{proof}
一个应用如下。取有向 $\Ai$ 范畴 $\A^{\to}$，以及 $\A=\Tw\A^{\to}$ 中由 $\A^{\to}$ 本身的对象构成的满例外集。按\secref{5j}施加变异来变换这个例外集，令 $\B^{\to}$ 为 $\A$ 中与变异后对象族相应的有向 $\Ai$ 子范畴。\stref{引理}{5.22}给出拟等价 $\Tw\B^{\to}\to\Tw\A^{\to}$，从而给出三角范畴的等价 $D(\B^{\to})\cong D(\A^{\to})$。正如 Kontsevich 最初观察到的，这定义了群 $\Z^m\rtimes\operatorname{Br}_m$ 在具有 $m$ 个对象的有向 $\Ai$ 范畴的拟同构类集合上的作用，同一轨道中的所有范畴都具有等价的导出范畴（实际上，中心元素 $(1,\ldots,1)\in\Z^m\subset\Z^m\rtimes\operatorname{Br}_m$ 作用平凡，因此这一作用下降到商群 $\Z^{m-1}\rtimes\operatorname{Br}_m$）。

我们还需要\stref{引理}{5.22}对非例外对象族的一个版本：
\begin{lemma}\label{st:5.23}
设 $Y$ 是 $\A$ 中的对象族，每个对象都是球面的。在 $\Tw\A$ 中，考虑由与\eref{5.24}相同的公式定义的对象族 $U=\mathscr L_kY$。用 $\B^{\to},\C^{\to}$ 表示与这两个对象族相应的有向 $\Ai$ 子范畴。那么 $\Tw\B^{\to}$ 与 $\Tw\C^{\to}$ 拟等价。
\end{lemma}
\begin{proof}
根据\stref{引理}{5.6}，可以把 $\Ai$ 范畴 $\A$ 替换为任意拟等价范畴，并把 $Y,U$ 替换为新范畴中的对应对象族。因此，可以假设 $\A$ 严格含幺且 $\mu^1_{\A}=0$；这纯粹是为了方便。

稍微区分一下记号，用 $\widetilde Y_k$ 表示 $\B^{\to}$ 中对应于 $\A$ 中 $Y_k$ 的对象。这些对象在 $\B=\Tw\B^{\to}$ 中构成满例外集 $\widetilde Y$。令 $\widetilde U=\mathscr L_k\widetilde Y$ 为这个例外集的第 $k$ 个基本变异。包含函子 $\mathcal H:\B^{\to}\to\A$ 诱导 $\Ai$ 函子 $\Tw\mathcal H:\B\to\Tw\A$，满足
\begin{equation}\label{eq:5.33}
\Tw\mathcal H(T_{\widetilde Y_i}\widetilde Y_j)=T_{Y_i}Y_j\quad(i<j).
\end{equation}
左边的扭转在 $\B$ 内进行，右边是在 $\Tw\A$ 中的相应构造。两边都采用\secref{5f}的显式公式，于是等式归结为 $\hom_{\B^{\to}}(\widetilde Y_i,\widetilde Y_j)=\hom_{\A}(Y_i,Y_j)$。由定义，\eref{5.33}尤其意味着 $\Tw\mathcal H(\widetilde U)=U$。下一步是证明映射
\begin{equation}\label{eq:5.34}
(\Tw\mathcal H)^1:\hom_{\B}(\widetilde U_i,\widetilde U_j)\longrightarrow\hom_{\Tw\A}(U_i,U_j)\quad(i<j)
\end{equation}
\SourcePage{92}
为拟同构。只有 $i$ 或 $j$ 等于 $k$ 时，这一点才非平凡。大多数情形可由以下观察处理：一边由 $p<q$ 时的 $\hom_{\B^{\to}}(\widetilde Y_p,\widetilde Y_q)$ 构成，另一边由相应的 $\hom_{\A}(Y_p,Y_q)$ 构成；依 $\B^{\to}$ 的定义，它们相同。唯一剩下的是 $i=k,j=k+1$ 的情形，此时\eref{5.34}中的空间为
\[
\begin{aligned}
\hom_{\B}(T_{\widetilde Y_k}\widetilde Y_{k+1},\widetilde Y_k)&=\hom_{\A}(Y_k,Y_{k+1})^{\vee}[-1],\\
\hom_{\Tw\A}(T_{Y_k}Y_{k+1},Y_k)&=\hom_{\A}(Y_{k+1},Y_k)\\[-2pt]
&\qquad\oplus\hom_{\A}(Y_k,Y_{k+1})^{\vee}[-1]\otimes\hom_{\A}(Y_k,Y_k).
\end{aligned}
\]
由于 $\mu^1_{\A}=0$，第一个空间的微分为零；第二个空间上的微分则是将复合部分对偶化所得的映射
\[
(\mu^2_{\A})':\hom_{\A}(Y_{k+1},Y_k)\longrightarrow\hom_{\A}(Y_k,Y_{k+1})^{\vee}\otimes\hom_{\A}(Y_k,Y_k).
\]
球面性意味着
\[
(\id\otimes\textstyle\int)\circ(\mu^2_{\A})':\hom_{\A}(Y_{k+1},Y_k)\longrightarrow\hom_{\A}(Y_k,Y_{k+1})^{\vee}[-n_k]
\]
是同构，其中 $n_k$ 为球面对象 $Y_k$ 的维数。因此 $(\mu^2_{\A})'$ 单射；再作维数计算，可知其像是与 $\ker(\id\otimes\int)=\hom_{\A}(Y_k,Y_{k+1})^{\vee}\otimes\K e_{Y_k}$ 互补的子空间。而\eref{5.34}中的映射恰好为 $(\Tw\mathcal H)^1=(0,\id\otimes e_{Y_k})$，因此是拟同构。

令 $\widetilde{\C}^{\to}$ 为 $\B$ 中与对象族 $\widetilde U$ 相应的有向 $\Ai$ 子范畴。它自然嵌入 $\B$；将此嵌入与 $\Tw\mathcal H$ 组合，得到 $\Ai$ 函子 $\G:\widetilde{\C}^{\to}\to\Tw\A$。这个函子实际上严格含幺，其像落在有向 $\Ai$ 子范畴 $\C^{\to}\subset\Tw\A$ 中；由上面的计算，它给出 $\widetilde{\C}^{\to}$ 与 $\C^{\to}$ 之间的拟同构。因此，$\Tw\C^{\to}$ 与 $\Tw\widetilde{\C}^{\to}$ 拟等价；而根据对变异的一般讨论，后一个范畴又与 $\Tw\B^{\to}$ 拟等价。
\end{proof}
要将有向子范畴的构造推广到 $\A$ 仅 c-有限的情形，可以这样做：利用扰动引理，转到一个拟同构的 $\Ai$ 范畴 $\widetilde{\A}$，使实际的 $\hom_{\widetilde{\A}}$ 空间有限维，再对这个范畴应用上述构造。将结果仍记为 $\A^{\to}$；它仍带有\stref{引理}{5.20}中的 $\Ai$ 函子 $\mathcal H:\A^{\to}\to\A$。缺点是 $\A^{\to}$ 不再典范，因为它依赖于 $\widetilde{\A}$ 的选择；但至少在拟同构意义下唯一，这对于大多数用途已经足够。特别地，\stref{引理}{5.22}与\stref{引理}{5.23}仍然成立。
